\section{Mô hình D-Tracker hoàn chỉnh}
\label{sec:model}

\subsection{Ba thành phần cộng lại}

Với cửa sổ hiện tại $\mathcal{X}^{c}$, paper giả định dữ liệu quan sát là xấp xỉ tổng của một
tensor trend, một tensor seasonal và một tensor outlier (Eq.~(1) trong paper):
\begin{equation}
  \mathcal{X}^{c}\;\approx\;\widehat{\mathcal{X}}_d
  +\widehat{\mathcal{X}}_s+\widehat{\mathcal{X}}_o,
  \qquad
  \widehat{\mathcal{X}}_d,\widehat{\mathcal{X}}_s,\widehat{\mathcal{X}}_o
  \in\mathbb{R}^{k\times l\times L_c}.
  \label{eq:model-sum}
\end{equation}
Mục đích của phép tách không phải là trang trí. Mỗi thành phần có một giả định riêng:
\begin{itemize}
  \item $\widehat{\mathcal{X}}_d$ chứa \textbf{trend và diffusion}, phần thay đổi có cấu trúc;
  \item $\widehat{\mathcal{X}}_s$ chứa \textbf{mùa vụ}, phần lặp lại theo chu kỳ;
  \item $\widehat{\mathcal{X}}_o$ chứa \textbf{outlier}, phần bất thường và được giả định thưa.
\end{itemize}
Nếu không tách, một đợt khuyến mãi có thể bị học như một xu hướng dài hạn; nếu không có outlier,
một cú sốc có thể làm hỏng mọi hệ số.

\subsection{Trend: nén lõi và chiếu trở lại}

Tensor trend vẫn có kích thước $k\times l\times L_c$, quá lớn để đọc trực tiếp. Paper phân rã nó
thành một lõi nhỏ và hai bảng hệ số (Eq.~(2) trong paper):
\begin{equation}
  \widehat{\mathcal{X}}_d
  =\mathcal{W}^{(\mathrm{core})}
   \times_{\mathrm{key}}\mathbf{W}^{(\mathrm{key})}
   \times_{\mathrm{loc}}\mathbf{W}^{(\mathrm{loc})},
  \label{eq:trend-factorization}
\end{equation}
với
\[
  \mathcal{W}^{(\mathrm{core})}\in\mathbb{R}^{d_k\times d_l\times L_c},
  \quad
  \mathbf{W}^{(\mathrm{key})}\in\mathbb{R}_{+}^{d_k\times k},
  \quad
  \mathbf{W}^{(\mathrm{loc})}\in\mathbb{R}_{+}^{d_l\times l},
\]
và $d_k\ll k$, $d_l\ll l$. Viết theo từng phần tử, công thức này là
\begin{equation}
  \widehat{x}_{d,uv}(t)
  =\sum_{i=1}^{d_k}\sum_{j=1}^{d_l}
    w^{(\mathrm{core})}_{ij}(t)\,
    W^{(\mathrm{key})}_{iu}W^{(\mathrm{loc})}_{jv}.
  \label{eq:trend-elementwise}
\end{equation}
Từng ký hiệu có nghĩa rất cụ thể: $u$ là từ khóa thật, $v$ là địa điểm thật, $i,j$ là nhóm ẩn.
Nếu $W^{(\mathrm{key})}_{iu}$ lớn, từ khóa $u$ ``thuộc về'' nhóm $i$; nếu
$W^{(\mathrm{loc})}_{jv}$ lớn, địa điểm $v$ thuộc nhóm $j$. Lõi $w^{(\mathrm{core})}_{ij}(t)$ là
đường hoạt động của cặp nhóm đó.

\begin{keypoint}
Đẳng thức~\eqref{eq:trend-elementwise} là trái tim của khả năng diễn giải: giá trị thật của một
cặp $(u,v)$ bằng tổng các dynamics ẩn, mỗi dynamics được nhân với mức liên hệ của từ khóa và địa
điểm. Hệ số không âm giữ cho phép cộng không đảo chiều.
\end{keypoint}

\subsection{Lõi được sinh bởi hệ reaction--diffusion}

Paper không cho $w^{(\mathrm{core})}_{ij}(t)$ biến đổi tự do. Nó giả định lõi được sinh bởi hệ
phương trình~\eqref{eq:rd-latent} ở Chương~\ref{sec:reaction-diffusion}, tức Eq.~(3) trong paper:
\begin{equation}
  \frac{dw_{ij}}{dt}
  =a_{ij}w_{ij}
   +\sum_{j'=1}^{d_l}d_{ijj'}(w_{ij'}-w_{ij}).
  \label{eq:core-ode}
\end{equation}
Đây là một hệ \textbf{phương trình vi phân thường (ODE)} theo thời gian. Dữ liệu được quan sát
theo các bước rời rạc (tuần hoặc ngày), nhưng mô hình định nghĩa động học liên tục theo $t$ rồi
tích phân để tạo ra chuỗi $L_c$ bước. Điểm cần nhớ: vị trí là rời rạc (các nhóm địa điểm), còn
thời gian trong phương trình là liên tục.

Gọi $\theta=\{\mathbf{A},\mathcal{D}\}$ là tham số của hệ, và $w_0\in\mathbb{R}_{+}^{d_k\times d_l}$
là trạng thái ban đầu. Paper ký hiệu $f_\theta(\theta,w_0,L)$ là hàm sinh ra một chuỗi lõi dài $L$
từ $w_0$. Như vậy, phần trend hoàn chỉnh có tham số
\[
  \theta_d=\{\mathbf{W}^{(\mathrm{key})},\mathbf{W}^{(\mathrm{loc})},\theta\}.
\]
Nói cách khác, $\theta_d$ chứa cả ``các dynamics tiến hóa thế nào'' và ``các dynamics được gắn
với từ khóa/địa điểm nào''.

\subsection{Seasonal: một phép phân rã tương tự}

Phần mùa vụ cũng được nén, nhưng không dùng ODE. Paper phân rã (Eq.~(4)):
\begin{equation}
  \widehat{\mathcal{X}}_s
  =\mathbf{S}^{(\mathrm{time})}
   \times_{\mathrm{key}}\mathbf{S}^{(\mathrm{key})}
   \times_{\mathrm{loc}}\mathbf{S}^{(\mathrm{loc})},
  \label{eq:seasonal-factorization}
\end{equation}
với $d_s$ pattern mùa vụ và
\[
  \mathbf{S}^{(\mathrm{time})}\in\mathbb{R}^{d_s\times L_c},
  \quad
  \mathbf{S}^{(\mathrm{key})}\in\mathbb{R}^{d_s\times k},
  \quad
  \mathbf{S}^{(\mathrm{loc})}\in\mathbb{R}^{d_s\times l}.
\]
Theo từng phần tử:
\begin{equation}
  \widehat{x}_{s,uv}(t)
  =\sum_{r=1}^{d_s}
    S^{(\mathrm{time})}_{r,t}S^{(\mathrm{key})}_{r,u}S^{(\mathrm{loc})}_{r,v}.
  \label{eq:seasonal-elementwise}
\end{equation}
Mỗi $r$ là một ``khuôn'' mùa vụ theo thời gian; hai hệ số còn lại nói khuôn đó mạnh thế nào ở từng
từ khóa và địa điểm. Đây là dạng phân rã PARAFAC: mỗi pattern là một tổng hạng nhất (rank-one) của
ba vector. Paper không nói các hệ số mùa vụ bắt buộc không âm như $\mathbf{W}$ của trend, nên tài
liệu này không tự thêm ràng buộc đó.

\subsection{Outlier: thưa, không phải mọi nhiễu}

Tensor outlier $\widehat{\mathcal{X}}_o$ được giả định \emph{sparse}: hầu hết phần tử bằng 0, chỉ
giữ những giá trị bất thường không giải thích được bởi trend và mùa vụ. Trong tối ưu, paper tách
residual $\mathcal{X}^{c}-\widehat{\mathcal{X}}_d-\widehat{\mathcal{X}}_s$ và chỉ giữ các phần tử
làm giảm tổng chi phí MDL. Vì vậy ``outlier'' ở đây là một thành phần có cấu trúc, không phải
nhãn ``dữ liệu xấu''.

\subsection{Bộ tham số và nhiều regime}

Một mô hình đơn lẻ được gói trong
\[
  \Theta=\{\theta_d,\theta_s,\widehat{\mathcal{X}}_o\},
  \qquad
  \theta_s=\{\mathbf{S}^{(\mathrm{time})},\mathbf{S}^{(\mathrm{key})},
             \mathbf{S}^{(\mathrm{loc})}\}.
\]
Khi pattern thay đổi, một $\Theta$ không còn đủ. Paper cho phép dùng $N$ mô hình:
\[
  \mathcal{F}=\{\Theta_1,\Theta_2,\ldots,\Theta_N\}.
\]
Mỗi $\Theta_n$ mô tả một giai đoạn tương đối ổn định; $N$ và thời điểm chuyển được chọn tự động
bằng MDL ở Chương~\ref{sec:mdl-streaming}, không phải chia đoạn bằng mắt. Vì vậy, khi đọc một
hình như Hình~\ref{fig:preview}(a), các đường chấm dọc không phải sự kiện dữ liệu; chúng là ranh
giới của các mô hình.

\begin{center}
\small
\begin{tabular}{@{}p{0.24\linewidth}p{0.68\linewidth}@{}}
  \toprule
  Tham số & Chứa gì \\
  \midrule
  $\theta$ & $\mathbf{A}$ (growth) và $\mathcal{D}$ (diffusion) của hệ ẩn. \\
  $\theta_d$ & Lõi trend, hai bảng hệ số từ khóa/địa điểm và $\theta$. \\
  $\theta_s$ & Các pattern mùa vụ và hệ số từ khóa/địa điểm. \\
  $\Theta$ & Một mô hình đầy đủ: trend, mùa vụ và outlier. \\
  $\mathcal{F}$ & Tập nhiều $\Theta$, tương ứng các giai đoạn của stream. \\
  \bottomrule
\end{tabular}
\end{center}

\subsection{Đọc một dynamics từ lõi đến dữ liệu thật}

Giả sử mô hình học được $d_k=3$ nhóm từ khóa và $d_l=2$ nhóm địa điểm. Muốn hiểu ``Dynamics \#1
của Group 1'' trong Hình~\ref{fig:preview}(b), ta đi theo bốn bước:
\begin{enumerate}
  \item Xem đường lõi của dynamics đó đang tăng, giảm hay đổi hướng.
  \item Xem cột hệ số từ khóa tương ứng: từ khóa nào có trọng số cao? Trong ví dụ 2019--2021,
        paper cho thấy ``iPhone'' liên hệ mạnh với Dynamics \#1.
  \item Xem hàng hệ số địa điểm: nhóm nào có trọng số cao? Paper mô tả Group~1 gồm Hoa Kỳ, châu
        Âu và Trung Quốc; Group~2 là phần còn lại.
  \item Xem các $d_{ijj'}$ và dấu của $(w_{ij'}-w_{ij})$ để đọc hướng lan truyền. Paper mô tả
        Dynamics \#1 lan từ Group~1 sang Group~2, tức quan tâm ``iPhone'' lan từ các khu vực dẫn
        đầu ra phần còn lại của thế giới trong giai đoạn đó.
\end{enumerate}
Đây là \emph{diễn giải theo mô hình}, không phải nhãn do con người gán trước. Tên ``Group 1'' chỉ
có nghĩa sau khi học hệ số; nhóm không nhất thiết trùng với châu lục hay ngôn ngữ.

\begin{plainwords}
Muốn đọc kết quả, hãy tưởng tượng ba lớp: lõi nhỏ nói ``chuyện gì đang xảy ra'', hàng/cột hệ số
nói ``ở đâu và với từ khóa nào'', còn mũi tên diffusion nói ``đi từ nhóm nào sang nhóm nào''.
Bạn không cần nhìn toàn bộ tensor gốc.
\end{plainwords}

\begin{commonmistake}
Một lỗi phổ biến là đọc $d_l\ll k$ trong phần mô tả của paper như một ràng buộc thật. Kích thước
của $\mathbf{W}^{(\mathrm{loc})}$ là $d_l\times l$ và ý tưởng là giảm số nhóm địa điểm, nên tài
liệu dùng $d_k\ll k$, $d_l\ll l$. Đây là chỗ paper có khả năng đánh máy; ta nêu rõ thay vì lặng lẽ
sao chép. Một lỗi khác là coi $\Theta$ và $\mathcal{F}$ là một; $\mathcal{F}$ là \emph{tập} các
mô hình cho các giai đoạn khác nhau.
\end{commonmistake}

\begin{selfcheck}
\begin{enumerate}
  \item Viết kích thước của bốn đối tượng trong Công thức~\eqref{eq:trend-factorization}.
  \item Phân biệt $\theta$, $\theta_d$, $\Theta$ và $\mathcal{F}$ trong một câu mỗi ký hiệu.
  \item Vì sao seasonal không cần ODE reaction--diffusion? Nó khác trend ở giả định nào?
  \item Nếu $N=1$ trong toàn bộ stream, mô hình đang giả định điều gì?
\end{enumerate}
\end{selfcheck}
